\subsection{Connecting homomorphisms and exact sequences}

Let M be a right A-module. Recall that for every left A-module N, we defined in the preceding number (X, p.~69, prop. 5) an isomorphism

Let

\[
\psi_ {\mathbf {M}} (\mathbf {N}): \operatorname{Tor} ^ {\mathbf {A}} (\mathbf {M}, \mathbf {N}) \rightarrow \mathrm{H} (\mathrm{L} (\mathbf {M}) \otimes_ {\mathbf {A}} \mathbf {N}).\tag{ε}
\]

\[
0 \rightarrow \mathbf {N} ^ {\prime} \xrightarrow {u} \mathbf {N} \xrightarrow {v} \mathbf {N} ^ {\prime \prime} \rightarrow 0
\]

an exact sequence of left A-modules; the sequence of \(k\) -complexes

\[
\left(^ {\mathrm{M}} \mathcal {E}\right) \quad 0 \longrightarrow \mathrm{L} (\mathrm{M}) \otimes_ {\mathrm{A}} \mathrm{N} ^ {\prime} \xrightarrow {1 \otimes u} \mathrm{L} (\mathrm{M}) \otimes_ {\mathrm{A}} \mathrm{N} \xrightarrow {1 \otimes v} \mathrm{L} (\mathrm{M}) \otimes_ {\mathrm{A}} \mathrm{N} ^ {\prime \prime} \longrightarrow 0
\]

is then exact (X, p.~66, lemma 1); let

\[
\hat {c} \left(^ {\mathrm{M}} \mathcal {E}\right): \mathrm{H} (\mathrm{L} (\mathrm{M}) \otimes_ {\mathrm{A}} \mathrm{N} ^ {\prime \prime}) \rightarrow \mathrm{H} (\mathrm{L} (\mathrm{M}) \otimes_ {\mathrm{A}} \mathrm{N} ^ {\prime})
\]

the corresponding connecting homomorphism (X, p.~29).

DEFINITION 2. --- A connecting homomorphism of torsion products, relative to the module M and to the exact sequence \(\mathcal{E}\) , the composite homomorphism

\[
\partial (\mathrm{M}, \mathcal {E}) = \psi_ {\mathrm{M}} \left(\mathrm{N} ^ {\prime}\right) ^ {- 1} \circ \partial \left(^ {\mathrm{M}} \mathcal {E}\right) \circ \psi_ {\mathrm{M}} \left(\mathrm{N} ^ {\prime \prime}\right): \operatorname{Tor} ^ {\mathrm{A}} \left(\mathrm{M}, \mathrm{N} ^ {\prime \prime}\right)\rightarrow \operatorname{Tor} ^ {\mathrm{A}} \left(\mathrm{M}, \mathrm{N} ^ {\prime}\right).
\]

This is a graded \(k\) -homomorphism of degree \((-1)\) , whose homogeneous components are denoted \(\partial_n(\mathbf{M}, \mathcal{E}) : \mathrm{Tor}_n^{\mathbf{A}}(\mathbf{M}, \mathbf{N}'') \to \mathrm{Tor}_{n-1}^{\mathbf{A}}(\mathbf{M}, \mathbf{N}')\) .

THEOREM 1. --- The left-unbounded sequence of homomorphisms of k-modules

\[
\begin{array}{l} \dots \longrightarrow \operatorname{Tor} _ {n} ^ {\mathbf {A}} (\mathbf {M}, \mathbf {N} ^ {\prime}) \xrightarrow {\operatorname{Tor} _ {n} ^ {\mathbf {A}} (1 , u)} \operatorname{Tor} _ {n} ^ {\mathbf {A}} (\mathbf {M}, \mathbf {N}) \xrightarrow {\operatorname{Tor} _ {n} ^ {\mathbf {A}} (1 , v)} \operatorname{Tor} _ {n} ^ {\mathbf {A}} (\mathbf {M}, \mathbf {N} ^ {\prime \prime}) \\ \xrightarrow {\hat {c} _ {n} (\mathbf {M} , \mathcal {E})} \operatorname{Tor} _ {n - 1} ^ {\mathbf {A}} (\mathbf {M}, \mathbf {N} ^ {\prime}) \xrightarrow {\operatorname{Tor} _ {n - 1} ^ {\mathbf {A}} (1 , u)} \dots \xrightarrow {\operatorname{Tor} _ {1} ^ {\mathbf {A}} (1 , v)} \operatorname{Tor} _ {1} ^ {\mathbf {A}} (\mathbf {M}, \mathbf {N} ^ {\prime \prime}) \\ \xrightarrow {\hat {c} _ {1} (\mathbf {M} , \mathcal {E})} \mathbf {M} \otimes_ {\mathbf {A}} \mathbf {N} ^ {\prime} \xrightarrow {1 \otimes u} \mathbf {M} \otimes_ {\mathbf {A}} \mathbf {N} \xrightarrow {1 \otimes v} \mathbf {M} \otimes_ {\mathbf {A}} \mathbf {N} ^ {\prime \prime} \longrightarrow 0 \end{array}
\]

is exact.

Consider indeed the diagram

\[
\begin{array}{c} \operatorname{Tor} (\mathbf {M}, \mathbf {N} ^ {\prime}) \xrightarrow {\operatorname{Tor} (1 , u)} \operatorname{Tor} (\mathbf {M}, \mathbf {N}) \xrightarrow {\operatorname{Tor} (1 , v)} \operatorname{Tor} (\mathbf {M}, \mathbf {N} ^ {\prime \prime}) \xrightarrow {\partial (\mathbf {M} , \mathcal {E})} \operatorname{Tor} (\mathbf {M}, \mathbf {N} ^ {\prime}) \xrightarrow {\operatorname{Tor} (1 , u)} \operatorname{Tor} (\mathbf {M}, \mathbf {N}) \\ \psi_ {\mathbf {M} (\mathbf {N} ^ {\prime})} \Big | _ {\downarrow} \quad \psi_ {\mathbf {M} (\mathbf {N})} \Big | _ {\downarrow} \quad \psi_ {\mathbf {M} (\mathbf {N} ^ {\prime \prime})} \Big | _ {\downarrow} \quad \psi_ {\mathbf {M} (\mathbf {N} ^ {\prime})} \Big | _ {\downarrow} \quad \psi_ {\mathbf {M} (\mathbf {N} ^ {\prime})} \Big | _ {\downarrow} \\ H (L (\mathbf {M}) \otimes N ^ {\prime}) \xrightarrow {H (1 \otimes u)} H (L (\mathbf {M}) \otimes N) \xrightarrow {H (1 \otimes v)} H (L (\mathbf {M} \otimes N ^ {\prime \prime}) \xrightarrow {\partial^ {(M)} \mathcal {E}} H (L (\mathbf {M}) \otimes N ^ {\prime}) \xrightarrow {} H (L (\mathbf {M}) \otimes N). \end{array}
\]

It is commutative by (X, p.~69, remark 3) and def. 2. On the other hand, the bottom row is exact (X, p.~30, th. 1), and the various \(\psi_{\mathrm{M}}\) are bijective (X, p.~69, prop. 5).

COROLLARY 1. --- If Tor \(_{1}^{A}\) (M, N'') = 0, the sequence

\[
0 \longrightarrow \mathrm{M} \otimes_ {\mathrm{A}} \mathrm{N} ^ {\prime} \xrightarrow {1 \otimes u} \mathrm{M} \otimes_ {\mathrm{A}} \mathrm{N} \xrightarrow {1 \otimes v} \mathrm{M} \otimes_ {\mathrm{A}} \mathrm{N} ^ {\prime \prime} \longrightarrow 0
\]

is exact.

COROLLARY 2. --- Let \(0 \to C' \xrightarrow{u} C \xrightarrow{v} C'' \to 0\) be an exact sequence of complexes of left A-modules and E a complex of right A-modules. If \(C''\) or E is flat, the sequence

\[
0 \longrightarrow \mathrm{E} \otimes_ {\mathrm{A}} \mathrm{C} ^ {\prime} \xrightarrow {1 \otimes u} \mathrm{E} \otimes_ {\mathrm{A}} \mathrm{C} \xrightarrow {1 \otimes v} \mathrm{E} \otimes_ {\mathrm{A}} \mathrm{C} ^ {\prime \prime} \longrightarrow 0
\]

is exact.

Indeed, \(\operatorname{Tor}_1^{\mathbf{A}}(\mathbf{E},\mathbf{C}^{\prime \prime}) = 0\) according to X, p.~69, cor. to prop. 5.

Example. --- Let a be an ideal of A. The exact sequence

\[
0 \rightarrow \mathfrak {a} \rightarrow \mathrm{A} _ {\mathrm{s}} \rightarrow \mathrm{A} / \mathfrak {a} \rightarrow 0
\]

of left A-modules gives rise to an exact sequence of torsion products, in which the terms \(\operatorname{Tor}_{i}^{\mathbf{A}}(\mathbf{M},\mathbf{A})\) are zero for i \textgreater{} 0. From this, we deduce isomorphisms

\[
\operatorname{Tor} _ {i + 1} ^ {\mathbf {A}} (\mathbf {M}, \mathbf {A} / \mathfrak {a}) \rightarrow \operatorname{Tor} _ {i} ^ {\mathbf {A}} (\mathbf {M}, \mathfrak {a}), \quad i > 0
\]

and an exact sequence

\[
0 \rightarrow \operatorname{Tor} _ {1} ^ {\mathbf {A}} (\mathrm{M}, \mathrm{A} / \mathfrak {a}) \rightarrow \mathrm{M} \otimes_ {\mathbf {A}} \mathfrak {a} \rightarrow \mathrm{M} \otimes_ {\mathbf {A}} \mathrm{A} \rightarrow \mathrm{M} \otimes \mathrm{A} / \mathfrak {a} \rightarrow 0:
\]

it follows that \(\operatorname{Tor}_{1}^{\mathbf{A}}(\mathbf{M},\mathbf{A}/\mathfrak{a})\) is identified with the kernel of the canonical homomorphism \(\mathbf{M}\otimes_{\mathbf{A}}\mathfrak{a}\to\mathbf{M}\) .

For example, taking for M a module of the form \(A_{d}/b\) where b is a right ideal of A, one obtains an isomorphism of \(\operatorname{Tor}_{1}^{A}(A/b, A/a)\) onto \((a \cap b)/ba\) .

PROPOSITION 9. --- Let \(f: \mathbf{M} \to \mathbf{M}_{1}\) be a homomorphism of right A-modules and

(ε)

\[
\begin{array}{c} 0 \longrightarrow \mathrm{N} ^ {\prime} \stackrel {{u}} {{\longrightarrow}} \mathrm{N} \stackrel {{v}} {{\longrightarrow}} \mathrm{N} ^ {\prime \prime} \longrightarrow 0 \\ g ^ {\prime} \Biggl \downarrow \qquad \qquad \qquad \qquad \qquad \qquad \qquad \qquad \qquad \qquad \qquad \qquad \qquad \qquad \qquad \qquad \qquad \qquad \qquad \qquad \qquad \qquad \qquad \qquad \qquad \qquad \qquad \qquad \qquad \qquad \qquad \qquad \qquad \qquad \\ 0 \longrightarrow \mathrm{N} _ {1} ^ {\prime} \stackrel {{u _ {1}}} {{\longrightarrow}} \mathrm{N} _ {1} \stackrel {{v _ {1}}} {{\longrightarrow}} \mathrm{N} _ {1} ^ {\prime \prime} \longrightarrow 0 \end{array}\tag{$(\mathcal{E}_1)$}
\]

a commutative diagram with exact rows of homomorphisms of left A-modules. The diagram of k-modules

\[
\begin{array}{c} \operatorname{Tor} ^ {\mathbf {A}} (\mathbf {M}, \mathbf {N} ^ {\prime \prime}) \xrightarrow {\hat {c} (\mathbf {M} , \mathcal {E})} \operatorname{Tor} ^ {\mathbf {A}} (\mathbf {M}, \mathbf {N} ^ {\prime}) \\ \operatorname{Tor} ^ {\mathbf {A}} (f, g ^ {\prime \prime}) \Big \downarrow \\ \operatorname{Tor} ^ {\mathbf {A}} (\mathbf {M} _ {1}, \mathbf {N} _ {1} ^ {\prime \prime}) \xrightarrow {\hat {c} (\mathbf {M} _ {1} , \mathcal {E} _ {1})} \operatorname{Tor} ^ {\mathbf {A}} (\mathbf {M} _ {1}, \mathbf {N} _ {1} ^ {\prime}) \end{array}
\]

is commutative.

This follows from X, p.~31, prop. 2, applied to the commutative diagram

\[
\begin{array}{l} 0 \longrightarrow \mathrm{L(M)} \otimes_ {\mathrm{A}} \mathrm{N} ^ {\prime} \xrightarrow {1 \otimes u} \mathrm{L(M)} \otimes_ {\mathrm{A}} \mathrm{N} \xrightarrow {1 \otimes v} \mathrm{L(M)} \otimes_ {\mathrm{A}} \mathrm{N} ^ {\prime \prime} \longrightarrow 0 \\ \quad \mathrm{L(f)} \otimes g ^ {\prime} \Biggl \downarrow \quad \mathrm{L(f)} \otimes g \Biggl \downarrow \quad \mathrm{L(f)} \otimes g ^ {\prime \prime} \Biggl \downarrow \\ 0 \longrightarrow \mathrm{L(M} _ {1}) \otimes_ {\mathrm{A}} \mathrm{N} _ {1} ^ {\prime} \xrightarrow {1 \otimes u _ {1}} \mathrm{L(M} _ {1}) \otimes_ {\mathrm{A}} \mathrm{N} _ {1} \xrightarrow {1 \otimes v _ {1}} \mathrm{L(M} _ {1}) \otimes_ {\mathrm{A}} \mathrm{N} _ {1} ^ {\prime \prime} \longrightarrow 0. \end{array}
\]

Analogously, if N is a left A-module and

\[
0 \to \mathbf {M} ^ {\prime} \xrightarrow {r} \mathbf {M} \xrightarrow {s} \mathbf {M} ^ {\prime \prime} \to 0\tag{F}
\]

an exact sequence of right A-modules, one defines connecting homomorphisms

\[
\partial (\mathcal {F}, \mathrm{N}): \operatorname{Tor} ^ {\mathrm{A}} \left(\mathrm{M} ^ {\prime \prime}, \mathrm{N}\right)\rightarrow \operatorname{Tor} ^ {\mathrm{A}} \left(\mathrm{M} ^ {\prime}, \mathrm{N}\right)
\]

\[
\partial_ {n} (\mathcal {F}, \mathrm{N}): \operatorname{Tor} _ {n} ^ {\mathrm{A}} \left(\mathrm{M} ^ {\prime \prime}, \mathrm{N}\right)\rightarrow \operatorname{Tor} _ {n - 1} ^ {\mathrm{A}} \left(\mathrm{M} ^ {\prime}, \mathrm{N}\right)
\]

by \(\partial(\mathcal{F},\mathrm{N})=\overline{\psi}_{\mathrm{N}}(\mathrm{M}^{\prime})^{-1}\circ\partial(\mathcal{F}^{\mathrm{N}})\circ\overline{\psi}_{\mathrm{N}}(\mathrm{M}^{\prime\prime})\) , where \(\partial(\mathcal{F}^{\mathrm{N}})\) is the connecting homomorphism of the exact sequence

\[
(\mathcal {F} ^ {N})
\]

\[
0 \rightarrow \mathrm{M} ^ {\prime} \otimes_ {\mathrm{A}} \mathrm{L} (\mathrm{N}) \rightarrow \mathrm{M} \otimes_ {\mathrm{A}} \mathrm{L} (\mathrm{N}) \rightarrow \mathrm{M} ^ {\prime \prime} \otimes_ {\mathrm{A}} \mathrm{L} (\mathrm{N}) \rightarrow 0
\]

deduced from \(\mathcal{F}\) , and we have:

THEOREM 1 bis. --- The left-unbounded sequence of homomorphisms of k-modules

\[
\begin{array}{l}\rightarrow \operatorname{Tor} _ {n} ^ {\mathsf {A}} (\mathbf {M} ^ {\prime}, \mathbf {N}) \xrightarrow {\operatorname{Tor} _ {n} ^ {\mathsf {A}} (r , 1)} \operatorname{Tor} _ {n} ^ {\mathsf {A}} (\mathbf {M}, \mathbf {N}) \xrightarrow {\operatorname{Tor} _ {n} ^ {\mathsf {A}} (s , 1)} \operatorname{Tor} _ {n} ^ {\mathsf {A}} (\mathbf {M} ^ {\prime \prime}, \mathbf {N}) \xrightarrow {\partial_ {n} (\mathcal {F} , \mathbf {N})} \operatorname{Tor} _ {n - 1} ^ {\mathsf {A}} (\mathbf {M} ^ {\prime}, \mathbf {N})\\\dots \rightarrow \operatorname{Tor} _ {1} ^ {\mathsf {A}} (\mathbf {M} ^ {\prime \prime}, \mathbf {N}) \xrightarrow {\partial_ {1} (\mathcal {F} , \mathbf {N})} \mathbf {M} ^ {\prime} \otimes_ {\mathbf {A}} \mathbf {N} \xrightarrow {r \otimes 1} \mathbf {M} \otimes_ {\mathbf {A}} \mathbf {N} \xrightarrow {s \otimes 1} \mathbf {M} ^ {\prime \prime} \otimes_ {\mathbf {A}} \mathbf {N} \longrightarrow 0\end{array}
\]

is exact.

We leave it to the reader to state and prove the properties analogous to the corollaries of Thm. 1 and to Prop. 9. Moreover:

PROPOSITION 10. --- Let ( \(F^{\circ}\) ) denote the exact sequence of left A-modules

\[
0 \rightarrow \mathrm{M} ^ {\prime \circ} \xrightarrow {r} \mathrm{M} ^ {\circ} \xrightarrow {s} \mathrm{M} ^ {\prime \prime \circ} \rightarrow 0.
\]

The diagram

\[
\begin{array}{c} \operatorname{Tor} ^ {\mathbf {A}} (\mathbf {M} ^ {\prime \prime}, \mathbf {N}) \xrightarrow {\partial (\mathcal {F} , \mathbf {N})} \operatorname{Tor} ^ {\mathbf {A}} (\mathbf {M} ^ {\prime}, \mathbf {N}) \\ \sigma_ {\mathbf {M} ^ {\prime}, \mathbf {N}} \Biggl \downarrow \qquad \qquad \qquad \qquad \qquad \qquad \qquad \qquad \qquad \sigma_ {\mathbf {M} ^ {\prime}, \mathbf {N}} \Biggl \downarrow \\ \operatorname{Tor} ^ {\mathbf {A} ^ {\circ}} (\mathbf {N} ^ {\circ}, \mathbf {M} ^ {\prime \prime^ {\circ}}) \xrightarrow {\partial (\mathbf {N} ^ {\circ} , \mathcal {F} ^ {\circ})} \operatorname{Tor} ^ {\mathbf {A} ^ {\circ}} (\mathbf {N} ^ {\circ}, \mathbf {M} ^ {\prime \circ}) \end{array}
\]

is commutative.

Indeed, this follows from X, p.~31, prop. 2, applied to the commutative diagram

\[
\begin{array}{l} 0 \longrightarrow \mathrm{M} ^ {\prime} \otimes_ {\mathrm{A}} \mathrm{L} (\mathrm{N}) \xrightarrow {r \otimes 1} \mathrm{M} \otimes_ {\mathrm{A}} \mathrm{L} (\mathrm{N}) \xrightarrow {s \otimes 1} \mathrm{M} ^ {\prime \prime} \otimes_ {\mathrm{A}} \mathrm{L} (\mathrm{N}) \longrightarrow 0 \\ \sigma (\mathrm{M} ^ {\prime}, \mathrm{L} (\mathrm{N})) \Biggl \downarrow \quad \sigma (\mathrm{M}, \mathrm{L} (\mathrm{N})) \Biggl \downarrow \quad \sigma (\mathrm{M} ^ {\prime \prime}, \mathrm{L} (\mathrm{N})) \Biggl \downarrow \\ 0 \longrightarrow \mathrm{L} (\mathrm{N} ^ {\circ}) \otimes_ {\mathrm{A} ^ {\circ}} \mathrm{M} ^ {\prime \circ} \xrightarrow {1 \otimes r} \mathrm{L} (\mathrm{N} ^ {\circ}) \otimes_ {\mathrm{A} ^ {\circ}} \mathrm{M} ^ {\circ} \xrightarrow {1 \otimes s} \mathrm{L} (\mathrm{N} ^ {\circ}) \otimes_ {\mathrm{A} ^ {\circ}} \mathrm{M} ^ {\prime \prime \circ} \longrightarrow 0. \end{array}
\]

We shall see other commutation relations later (cf. X, p. 131, cor. 1).

